% !PRINTED 615
\weberpage{615}
\gesperrt{If $E$ denotes the group of units, or also a group contained in it, then we obtain the same value of the class number and the same class field if we let $EA$ take the place of $A$.}

This holds not only for this particular group, but generally. We could therefore, e.\,g., also have defined the numbers $\alpha$ of the group $A$ by
\[
\alpha \equiv \pm 1 \quad (\mathrm{mod}\ \mathfrak{m})
\]
since $E = +1, -1$ is a group also in the exceptional cases $\Delta = -4$, $\Delta = -3$.

Now it can be shown that the division field $\mathfrak{T}_{\mathfrak{m}}$ is the class field for $A$. If we first exclude all prime factors which divide the discriminant of the division field, and denote by $\tau$ any number of the division field, then by \S\,160
\weq{5}{\tau^p \equiv \tau \quad (\mathrm{mod}\ \mathfrak{p})}
is satisfied only when $\mathfrak{p}$ is a \gesperrt{prime ideal of the first degree} in $\Omega$ which belongs to the principal class. But that is the criterion for the class field. By \S\,154, however, the degree of the division field is at most equal to the expression $h$ given in (4); consequently it is exactly equal to $h$, and the division equation is irreducible. Moreover, it is thereby proved that in each of the ideal classes $\overline{A}_\nu$ with respect to $A$ there are infinitely many prime ideals of the first degree.

One can give this theorem the following expression:

If $\beta$ is an arbitrary number in $O$, and $\mathfrak{a}$ an arbitrary ideal in $\overline{O}$, then by
\weq{6}{\overline{A}_\nu = \mathfrak{a}\beta A}
a class with respect to $A$ is determined (\S\,161), and this class contains infinitely many prime ideals of the first degree. By (6) such a prime ideal can be represented in the form:
\weq{7}{\mathfrak{p} = \mathfrak{a}\beta\alpha,}
where $\alpha \equiv 1$ (mod $\mathfrak{m}$). If we therefore set
\[
\alpha\beta = \pi,
\]
then
\[
\pi \equiv \beta \quad (\mathrm{mod}\ \mathfrak{m}).
\]
From this it follows:

2. \gesperrt{If $\beta$ is a number in $O$, $\mathfrak{a}$ an ideal in $\overline{O}$, then there are infinitely many, in general fractional, numbers $\pi$ in $O$ which satisfy the condition}
%CONT>>
